\section{Introduction}

One in four preference battles hides a silent failure: reward models confident, humans conflicted. Reinforcement Learning from Human Feedback (RLHF) has become the dominant paradigm for aligning large language models with human preferences. Central to this process are reward models trained to predict human judgments, evaluated on benchmarks like RewardBench that report aggregate accuracy metrics---often exceeding 85\%. Yet aggregate metrics conceal systematic failure patterns. When a reward model scores 85\% accuracy, what distinguishes the successful 85\% from the failed 15\%? More critically, within that successful majority, how many cases reflect genuine alignment versus confident misalignment---the model certain while humans disagree?

This paper introduces a decomposition framework that reveals structure hidden by aggregate evaluation. We classify preference battles along two orthogonal dimensions: \textbf{human vote entropy} (high when voters disagree, low when consensus) and \textbf{reward model variance} (high when the model is uncertain, low when confident). Crossing these dimensions yields four alignment modes:

\begin{itemize}
    \item \textbf{Mode~1 (Aligned-Confident)}: Low entropy, low variance---humans agree, RM agrees
    \item \textbf{Mode~2 (Aligned-Uncertain)}: Low entropy, high variance---humans agree, RM uncertain
    \item \textbf{Mode~3 (Misaligned-Confident)}: High entropy, low variance---humans disagree, RM confident
    \item \textbf{Mode~4 (Misaligned-Uncertain)}: High entropy, high variance---both disagree
\end{itemize}

Mode~3---overconfident misalignment---is the critical failure pattern. Here, reward models express high confidence on precisely the samples where human preferences diverge. If prevalent, this mode represents a systematic blind spot in RLHF training signal.

Analyzing 57,477 battles from the Chatbot Arena, we find Mode~3 constitutes \textbf{23.7\%} of samples (95\% CI: 23.4--24.0\%, $p<0.001$), far exceeding our 10\% significance threshold. This is not a fringe phenomenon but a substantial pattern affecting approximately one quarter of real-world preference data.

Having established the phenomenon's existence, we test two mechanistic hypotheses. First, we hypothesize that Mode~3 arises when responses are semantically divergent---similar enough for RMs to miss distinctions humans catch. Our analysis falsifies this: Mode~3 response pairs show \emph{higher} semantic similarity than Mode~1 pairs (Cohen's $d = -0.049$), the opposite direction predicted. Second, we hypothesize Mode~3 concentrates in subjective prompt types (creative writing) versus objective tasks (coding, math). The data shows only negligible enrichment (ratio $= 1.07$ vs.\ predicted 1.5); Mode~3 pervades all task types.

\textbf{Contributions.} (1) To our knowledge, we provide the first empirical quantification of overconfident RM misalignment in large-scale preference data, demonstrating that $\sim$24\% of battles exhibit this failure pattern. (2) We introduce the $2\times2$ mode decomposition (entropy $\times$ variance) as a diagnostic framework that reveals patterns aggregate metrics miss. (3) We rule out two candidate mechanisms---semantic divergence and prompt subjectivity---constraining the hypothesis space for future work.
